\chapter{Electric Charges, Fields, and Gauss's Law}

\section{Electric Charges} \label{sec:electric-charges}

\textbf{Electric charge} is denoted by the symbol \(q\) or sometimes \(Q\),
and its SI unit is the coulomb (C).

Charge can be either positive or negative,
with protons carrying a positive charge and electrons carrying a
negative charge.

The \textbf{law of conservation of  charge} states that charge can
not created or destroyed,
only transferred,
and the net charge has to remain constant.
There are two ways to charge and discharge (i.e., move charges):
\begin{itemize}
    \item Conduction -- friction (triboelectricity) or contact
    \item Induction -- contactless
\end{itemize}

\textbf{Conductors} allow electrons to flow freely,
while \textbf{insulators} do not.
\vfil

\section{Electric Force and Coulomb's Law}

Charged objects exert an \textbf{electric force} on other charged
objects around them.
The magnitude of the electric force around a charged particle or
sphere can be given by \textbf{Coulomb's law},
which states
\begin{equation} \label{eq:coulombs-law}
    \begin{aligned}
        F &= k \frac{q_1 q_2}{r^2} \\
        &= \frac{q_1 q_2}{4 \pi \epsilon_0 r^2},
    \end{aligned}
\end{equation}
where
\begin{itemize}
    \item \(F\) is the force that \emph{both} objects experience.
    \item \(k = \frac{1}{4 \pi \epsilon_0}\) is the \textbf{Coulomb constant}.
    \item \(q_1\) and \(q_2\) are the two charges of the two objects.
        Note that both values are signed.
    \item \(r\) is the distance between the two particles.
    \item \(\epsilon_0 = \qty{8.854 187 8188(14)e-12}{\farad \per \metre}\)
        is the \textbf{permittivity of free space}
        or \textbf{vacuum permittivity}. \autocite{nist-epsilon_0}
\end{itemize}
Objects with charges of the same sign will repel each other,
while objects with charges of opposite signs will attract each other.
Thus, Coulomb's law will give a \emph{negative} value for attraction
and a \emph{positive} value for repulsion.

\section{Electric Fields}

\textbf{Electric fields} are closely related to electric forces
and have SI units of newtons per coulomb (\unit{\newton \per \coulomb}).
They are useful for cases where we only know the charge of one particle,
or wish to find the force exerted on multiple different particles.

\begin{subequations}
    The magnitude of the a charged sphere's electric field is given by
    \begin{equation} \label{eq:sphere-electric-field}
        \begin{aligned}
            E &= k \frac{q}{r^2} \\
            &= \frac{q}{4 \pi \epsilon_0 r^2}.
        \end{aligned}
    \end{equation}
    In vector form, the electric field is
    \begin{equation} \label{eq:sphere-electric-field-vector}
        \begin{aligned}
            \vecb{E} &= k \frac{q}{\|\vecb{r}\|^2} \hat{\vecb{r}} \\
            &= \frac{q}{4 \pi \epsilon_0 \|\vecb{r}\|^2}  \hat{\vecb{r}},
        \end{aligned}
    \end{equation}
    where \(\vecb{r}\) is the difference in position between the
    sphere exerting the field and the particle receiving it,
    i.e., \(\vecb{r} = \vecb{r_\mathrm{receive}} - \vecb{r_\mathrm{exert}}\).
\end{subequations}
Positive charges experience a force in the direction of the field,
and negative charges experience a force in the opposite direction.

To find the force exerted on a charged particle,
we multiply the electric field it is in by its charge,
i.e.,
\begin{equation}
    \vecb{F_E} = \vecb{E} q,
\end{equation}
where \(\vecb{F_E}\) is the electric force,
\(\vecb{E}\) is the electric field,
and \(q\) is the charge.

\subsection{Vector Fields}

We can draw electric fields as a vector field,
with arrrows representing the direction of the field
and their size representing the magnitude.

% To-do: add drawing of vector fields.

\subsection{Uniform Electric Fields}

The line between a positive and negative charge,
regardless of their respective magnitudes,
is a \textbf{uniform electric field}.
Outside of this line,
the electric field is not uniform due to \textbf{edge effects}.
This idea comes back up again in \pgRef{sec:capacitors}.

\section{Electric Flux and Gauss's Law}

\textbf{Flux}, symbolized by \(\Phi\),
is the amount of field that passes through a \textbf{surface},
which may be imaginary.
The word "flux" is derived from the Latin word for "flow"
in analogy with flowing water.

The general formula for flux is
\begin{equation}
    \Phi = \iint_S \vecb{F}(\vecb{r}) \cdot \odif{\vecb{S}},
\end{equation}
where \(S\) is the surface to integrate over
and \(\vecb{F}(\vecb{r})\) is the function of the field in 3D space.

For introductory electricity and magnetism,
we usually focus on a few simpler cases.

Surfaces can be either open (e.g., a window pane) or closed (e.g., a box).

For the case of a uniform electric field passing through an flat open surface,
the flux equation simplifies to
\begin{equation}
    \begin{aligned}
        \Phi_E = \vecb{E} \cdot \vecb{A} \\
        \Phi_E = E A \cos\theta
    \end{aligned}
\end{equation}
where
\begin{itemize}
    \item \(\Phi_E\) is the electric flux,
    \item \(\vecb{E}\) is the electric field vector and \(E\) is its magnitude,
    \item \(\vecb{A}\) is area vector and \(A\) is its magnitude,
        Note \(\vecb{A}\) is in the direction of the area's normal vector,
    \item \(\theta\) is the angle between the electric field
        and the area's normal vector.
\end{itemize}
Our surface in this case is infinitely thin.
By intuition, when \(\theta\) is \(\qty{90}{\degree}\),
there is no area to intercept the electric field,
thus the flux is 0.

The case of a closed surface is slightly more complicated.
For a closed surface, fields can flow both in and out.
When a field enters a closed surface, the flux is negative,
and when a field exits, the flux is positive.
As a result,
all fields that originate outside the closed surface result in a net flux of 0.

This means that only electric charges inside a closed surface
contributes to the electric flux,
which leads us to \textbf{Gauss's law} in integral form:
\begin{equation} \label{eq:gausss-law}
    \oiint \vecb{E} \cdot \odif{\vecb{A}} = \frac{q}{\epsilon_0},
\end{equation}
where
\begin{itemize}
    \item \(\vecb{E}\) is the electric field,
    \item \(\vecb{A}\) is the area vector of the surface,
    \item \(q\) is the charge enclosed by the surface.
        Some authors may use \(q_\mathrm{en}\)
        to emphasized that it must be enclosed.
\end{itemize}

Gauss's law is a generalization of Coulomb's law,
so we can show that Gauss's law simplies down to Coulomb's law
for the case of a charge particle in a spheral surface.
Let
\begin{itemize}
    \item \(\vecb{A}\) be the area vector of the spherical surface
        and \(A\) be its magnitude,
    \item \(r\) be the radius of the spherical surface,
    \item \(\vecb{E}\) be the electric field and \(E\) be its magnitude,
    \item \(\Phi_E\) is the electric flux,
    \item \(q\) be the enclosed charge.
\end{itemize}
The surface area of a sphere is given by
\begin{equation}
    A = 4 \pi r^2.
\end{equation}
Since electric fields radiate in a sphere,
the electric field and area are orthogonal at all points:
\begin{equation}
    \vecb{E} \cdot \odif{\vecb{A}} = E \odif{A}.
\end{equation}
If we integrate the electric field over the surface of the sphere, we find
\begin{equation}
    \Phi_E = 4 E \pi r^2.
\end{equation}
By Gauss's law, [\pgRef{eq:gausss-law}] we know that
\begin{equation}
    \begin{aligned}
        4 E \pi r^2 &= \frac{q}{\epsilon_0} \\
        E &= \frac{q}{4 \epsilon_0 \pi r^2}.
    \end{aligned}
\end{equation}
To get an electric force,
let \(q_1 = q\) and multiply by the other particle's charge \(q_2\):
\begin{equation*}
    \begin{aligned}
        4 E \pi r^2 &= \frac{q_1 q_2}{\epsilon_0} \\
        E &= \frac{q_1 q_2}{4 \epsilon_0 \pi r^2}.
    \end{aligned}
\end{equation*}
This leads us back to Coulomb's law [\pgRef{eq:coulombs-law}].
